\chapter{Payment and Transaction Issues}

\section{The Rush Is on to Make Net Commerce Happen}
\label{rush_commerce}

\subsection{Introduction}
The Internet is capable of becoming one of the most powerful business
and financial tools of our time--if the process becomes simple and
complete enough. Some businesses launched on the Net sell goods and
services, but the buyer must send a check via regular mail to complete
transactions. Although this is a viable method for some applications, a
richer suite of payment mechanisms is needed for the broader range of
business that will be conducted on the Internet.\\

Like the traditional off-line market, buyers and sellers want a variety
of options such as credit cards, debit cards, and cash--or in this
case, digital cash. To meet this demand, a host of new companies are
vying to establish their technologies as the means for conducting
secure Internet commerce.\\

Online commerce has existed in the form of EDI (electronic document
interchange) since the mid-1980s. EDI applications are
transaction-oriented, ranging from government RFPs (Requests for
Proposals) to purchase orders. EDI is in use throughout several major
industries to automate and simplify the otherwise paper-intensive
process of doing business. While performed today via secure Value-Added
Networks, it is questionable to what extent EDI applications will
translate to the Internet, especially since many of the transactions
involve large sums of money. A convergence is on the horizon
nevertheless.\\

\subsection{What is available? What is coming?}
CommerceNet, introduced in March 1994, is an industry consortium
organized into working groups. Its members comprise an intriguing mix
of Silicon Valley semiconductor manufacturers, software companies,
financial institutions, information providers, and businesses of
varying size. The consortium is a hothouse for electronic commerce and
a great source of information, especially for businesses. As of October
1994, CommerceNet had more than 50 members and 15 subscribers (small to
midsize companies).\\

Cisco Systems and Premenos Corp. recently announced a secure electronic
commerce experiment that should be readyin early 1995. The system uses
RSA Data Security Inc.'s public-key cryptography.\\

CyberCash Inc. is a well-funded start-up combining the talents of Dan
Lynch and William. Lynch converted the ARPAnet--the predecessor to the
Internet--to TCP/IP back in 1983 and has continually demonstrated his
abilities as an Internet visionary. Melton is the founder of Verifone,
a company that made real-time credit-card verification a reality for
retailers. Together they are working on a model for Internet commerce
that will accept digital cash, and that already deals with credit
cards, or debit cards. Expect to see CyberCash make its official debut
in the spring of 1995.\\

DigiCash was founded by Dr. David Chaum. The company pioneered E-Cash,
a software-only product that allows digital cash transfer over the
Internet. E-Cash was developed for smart cards and electronic
wallets--used in tollbridges, for example. DigiCash will supply the
technology through other organizations, including Encyclopedia
Britannica and NCSA (National Center for Supercomputing Applications),
which will release the products under various cooperation and trial
programs. The user software, which allows both paying and receiving
payment, will be distributed free of charge. DigiCash will encrypt
transactions and ensure anonymity.\\

First Virtual Holdings Inc. is attempting to become the first Internet
merchant banking system. The company brings together the Hollywood
business savvy of Lee Stein, First Virtual's president and CEO, and
three of the leading Internet gurus: Dr. Marshall T. Rose, Dr.
Nathaniel S. Borenstein, and Einar Stefferud.\\

Stein, an attorney and accountant, provides financial advisory services
to the music and entertainment industries. First Virtual principal
Rose, one of the Internet wunderkinds, has written the book (actually
several of them) on SNMP (Simple Network Management Protocol) and
network management. Borenstein, the chief scientist, is the primary
author of the MIME (Multipurpose Internet Mail Extensions) protocol
that allows transfer of anything digital over the Internet, including
multimedia documents. Stefferud, probably the world's foremost expert
on e-mail, has the role of chief visionary.\\

Introduced and online in October 1994, the First Virtual payment system
allows any Internet user to initiate retail transactions via e-mail.
The system aims to help maintain Internet culture by allowing buyers to
evaluate information prior to committing to a purchase, which is a
concept similar to shareware. No sensitive data is exchanged online in
First Virtual's system. Users set up accounts by telephone or fax, and
credit-card payments are processed off-line.\\

Registration fees are \$2 for consumers and \$10 for sellers. For each
transaction, sellers pay a 29-cent fee plus two percent of the
transaction price and a \$1 processing fee each time a payment is made
to their account. Initial First Virtual merchants include Internet
Resources Group, Internet Multicasting Service, and an Internet
newsgroup filtering service founded by David Farber, an Internet
luminary. The Department of Commerce is also considering adopting the
system to offset distribution costs of electronic documents.\\

The Internet Shopping Network was founded by its president, Randy
Adams, who formerly was CEO and chairman of IIAT Inc., a publicly held
international education and training organization. The company, which
was recently acquired by the Home Shopping Network, is an online
microcomputer software and hardware superstore available to anyone with
direct access to the Internet. Internet Shopping Network distributes
more than 20,000 computer software and hardware products, plus
InfoWorld, a computer news weekly. Membership is free, and prospective
members need only pre-register with an approved Visa or MasterCard.\\

Microsoft is also getting into the act. It recently announced a venture
with Visa International that will enable Windows users to make secure
credit-card purchases over the Internet using software co-developed by
the two companies. The software will be based on RSA Data Security's
encryption technology and Visa's VisaNet system will be used to
authenticate payments. According to Chairman Bill Gates `Fall Comdex
'94 speech, a `Wallet PC' that handles financial transactions and other
consumer retail interactions may also be part of our future.\\

Netscape Communications Corp., which provides an enhanced commercial
World-Wide Web browser, announced that it will support credit-card
purchases via the Internet through credit-card processor First Data
Corp. The Netscape system uses RSA's encryption technology to scramble
sensitive data. Merchants can purchase Netsite server software for
\$5,000, and buyers can access the system through free client software.\\

NetCash is another secure Internet project underway. Developed at the
Information Sciences Institute (ISI) of the University of Southern
California, it is a model for electronic currency that can be
anonymously transferred over an unsecure network without the use of
firewalls. NetCash allows service providers and users to select payment
mechanisms with varying degrees of anonymity.\\

Unlike most electronic cash systems in which only one financial
institution generates and accepts its own digital cash, NetCash is
designed to support the transfer of electronic cash between currency
servers. At this time, NetCash has not yet been made commercially
available.\\

Open Market is another secure Internet transactions venture. The system
supports secure payment, real-time credit-card authorization, account
statements, administrative interfaces for storefront management, and a
customer feedback mechanism. The company offers `document
fingerprinting, ` which automatically and uniquely numbers documents so
that fraudulent distribution can be tracked.\\

Open Market's StoreBuilder kit ranges in price from \$300 to \$1,500 for
setup and \$50 to \$300 in monthly fees, in addition to transaction fees
and additional storage fees. Buyers can set up an account by sending
credit-card info over the Net or by phone and fax. Open Market uses
password and encryption schemes to secure its accounts.\\

Finally, Terisa Systems is a joint venture of EIT (Enterprise
Integration Technologies) and RSA Data Security that is dedicated to
producing a secure Mosaic product suite for Internet commercial
transactions (see Section \ref{tool_commerce}).\\

\subsection{Connecting financial systems}
Electronic commerce players are teaming up with banks, credit-card, EDI,
and software companies, and more to create an infrastructure that can
support business over the Internet. EDS, Visa, and First USA --- which have
teamed up with First Virtual --- intend to stake a claim in the online
transaction business. In fact, participation from this type of company is
necessary so that people can buy goods and services in real time.\\

To make Internet commerce viable on a mass scale, it is necessary to
connect the world's financial systems. Given that security is
imperative to these organizations, why are they willing to risk doing
business on the Internet?\\

CyberCash and First Virtual are not directly connecting the Internet to
the global financial networks. CyberCash will use an intermediate
highly secure server that is connected to both networks. First
Virtual's system maintains security without having to use the
encryption or signature technologies that the other systems employ. The
system separates business on the Internet, which is insecure, from the
bank processing system through firewalls and off-line batch processing
handled by EDS, one of the world's leaders in banking transactions.


\subsection{No numbers, no crime}
Credit cards are currently the most popular way to buy something on the
Internet. This method is too dangerous to be practical, according to
First Virtual's Rose and Borenstein.\\

``It's very easy to write a program that looks for a sequence of
numbers beginning with a known credit-card prefix,'' Rose said. That is
why First Virtual never requires a potential information merchant or
consumer to send credit-card or checking-account numbers over the
Internet. Instead, First Virtual customers enter these numbers using an
automated phone system to set up a First Virtual account.\\

The Internet Shopping Network also keeps credit-card information off
the Net. Members enter their credit-card numbers using phone or fax
when they first establish an account. Open Market's system also
includes this option.\\

CyberCash deals with the problem using encryption. Its customers will
have a public key and a private key and must use both to complete a
transaction. The company will also use authentication and digital
signatures for added security.\\

Debit cards can be used for Internet-based financial transactions today
because the technology is in wide use among banks and retail points of
purchase. It is likely that the percentage of businesses accepting
debit cards for Internet purchases will be higher than the existing
number in the off-line marketplace because no special hardware is
required.\\

At the October World-Wide Web conference in Geneva, DigiCash CEO Chaum
demonstrated how anyone with a personal computer could transfer
DigiCash's E-Cash to any other workstation via e-mail on the Internet.
The E-Cash system was shown integrated with Mosaic.\\

Digital cash is the binary equivalent to currency. It differs from
credit and debit cards because it enables anonymous transactions.\\

``How much money are you willing to lose?'' asked CyberCash's Lynch.
``You don't, and I don't, walk around with \$10,000 in your wallet.
Usually people carry around what they are willing to lose. The same
will be true for digital cash.''\\

According to Chaum, DigiCash has a solution to the problem. ``E-Cash is
as secure as any government network because it employs the same
underlying security mechanisms,'' he said. E-Cash ensures the high
security required for electronic networks exclusively using `blind
signature' technology, an innovation in public-key cryptography.


\subsection{Internet-savvy business}
CyberCash is targeting banks and businesses, the Internet Shopping
Network and Open Market are targeting businesses and consumers, and
First Virtual is targeting all these and the government. In any case,
it appears that there will be a solution for just about everyone.\\

Over the past 18 months or so, organizations have been setting up WWW
and Gopher servers to distribute product and company information.
Businesses clever enough to have pointers strategically located to and
from other points on the Internet have taken the first steps in doing
real business on the Internet, but that is just the beginning.\\

For businesses, the Internet represents the cheapest new business
opportunity in history because many of the high costs of doing business ---
shipping, warehousing, duplication, and distribution --- are
handled at a fraction of the cost using servers and the network.\\

That said, businesses will not drop existing products and channels of
distribution just to go online. It is very conceivable, however, that
many businesses will use the Internet as an alternative form of
distribution because of its large reach and relatively low cost. Many
more small entrepreneurial businesses also will likely emerge if the
Internet proves to be a viable business alternative.


\section{A Tool to Make Web Commerce Secure}
\label{tool_commerce}

\subsection{Introduction}
With the explosion of interest in the World-Wide Web, hundreds, if not
thousands, of companies have rushed to get their home pages online to
provide general information, advertise, and offer catalogs. Many of
them would like to accept orders from customers online.\\

However, these firms generally have found the level of security in Internet
applications woefully inadequate. Companies want tools that allow them to
know with some certainty that customers are who they
claim to be so that billing can be charged against the proper accounts.\\

Similarly, consumers are concerned about security and privacy of online
transactions. They want to know that personal account information or
credit-card numbers will be seen only by authorized eyes. And as
vendors want to verify that a credit-card number is indeed a customer
`s, consumers want assurance that a person who purports to represent an
online service really belongs to that service.

\subsection{A head for business}
Several dozen companies are working on extensions to Internet protocols to
provide these types of security features (see Section \ref{rush_commerce}).
One of the first firms to announce such an effort was Silicon Valley's
Enterprise Integration Technologies (EIT).  EIT is working to provide
Internet tools that will allow retail firms and shoppers to conduct
business online as well as let businesses handle online transactions with
their suppliers.\\

One of EIT's endeavors is support of CommerceNet, a consortium of
California's Silicon Valley companies. The purpose of CommerceNet is to
develop Internet technologies that will secure commercial interaction
among firms in the Valley. A specific focus is to help foster their use
of the Internet for commerce among these firms and customers across the
globe.\\

Marty Tenenbaum, EIT's president, offered an example of an individual
customer who wants to buy a personal computer from a retailer in the
Valley. The customer would scan the company's online catalog using a
secure version of a World-Wide Web browser such as Mosaic; no special
security features would be required while the customer was browsing the
catalog. Should the customer decide to order a machine, he or she would
fill out a form with the desired options for the PC and submit the
purchase order to the Web server at the retailer.\\

This is where security features become important. Although many
consumers have sent credit-card numbers over the Internet via e-mail
and World-Wide Web forms, those who are knowledgeable about Internet
security are concerned about the information being intercepted, either
on local networks or on any of the intervening networks between the
consumer's desktop and the vendor's site. Tenenbaum's vision calls for
a secure Mosaic, which would use public-key cryptography to encrypt the
consumer's order to keep sensitive credit-card information safe from
prying eyes.

\subsection{The safe sell}
The use of a credit card is analogous to a consumer filling out an order
form and sending it in a paper envelope via postal mail. The customer is
reasonably certain that the order and credit-card number will arrive
safely. The vendor relies on a separate step --- authorization and credit
balance checking via the traditional credit-card
networks --- before fulfilling the order.\\

Once the PC vendor has received the order and verified the validity of
the customer's account, the vendor must ship the product. Many personal
computer vendors custom assemble each PC as orders arrive. In
Tenenbaum's plan, the PC vendor would also rely on the Internet as a
way for retail vendors to send orders to its suppliers.\\

Orders for everything from the chassis to the motherboard to the
monitor would be sent over the Internet. With computerized inventory
control, such orders could be placed online automatically, with the
retailer able to operate using just-in-time inventory methods.\\

In fact, Tenenbaum sees virtual corporations springing up as a result
of these technologies. Retail firms might become so intimately linked
to their suppliers--with order processing occurring so fluidly and
rapidly--that the nature of the relationship between suppliers and
retailers would be changedaltogether.\\

Encrypted credit-card transactions are only one type of business that
companies would like to conduct on the Internet. In the everyday
business world, signatures are an essential element--for signing
purchase orders, contracts, and the like. Most companies will accept
facsimile transmissions of some documents, but Internet e-mail is
easily enough `spoofed ` that its use is discouraged for these
applications. Some sort of digital signature is required to verify the
integrity of such documents.\\

The major elements of security required for commerce on the Internet,
therefore, are the following:
\begin{itemize}
  \item Authentication of individuals and companies so both parties can
    trust that the parties they are dealing with are who they claim to be.
  \item Encryption of digital documents so that sensitive information
    cannot be decoded even if the messages are intercepted.
  \item Digital signatures that verify the integrity of documents, so
    senders and recipients can trust that documents have not been altered
    in transit.
\end{itemize}

\subsection{Terisa is born}
To turn these desires into reality, EIT joined with RSA Data Security
Inc., a leading firm in digital-security technology, to form a company
called Terisa Systems. The venture's first goal is to create a suite of
client and server tools called the SecureWeb Toolkit. The tools will be
used by World-Wide Web client and server vendors to build products for
sale to end users. Terisa is working with the National Center for
Supercomputing Applications (NCSA) --- the organization that brought
Mosaic to the Net --- to build Secure NCSA Mosaic, which is expected to be
freely available on the Internet.\\

EIT and its partners have defined a standard called Secure Hypertext
Transfer Protocol, or Secure HTTP, for use in tools like Secure NCSA
Mosaic. The standard calls for extensions to the Web's basic HTTP
protocol to support needed security features Secure HTTP relies on
public-key encryption technology patented by RSA Data Security (see
Section \ref{data_crypto}).\\

How can users expect to cope with the encryption and authentication
process, which at first glance seems to be enormously complicated? In
theory, most of the details of the process are hidden to the end user.\\

A customer buys some sort of secure Mosaic package from a seller of
Internet products --- perhaps from his or her Internet service provider.
The package is preconfigured to work with a given public-key encryption
certifying authority (which verifies that the alleged owner of a public
key is genuine). The user installs the package on his or her computer,
and begins using it to browse the Web. When the user encounters a
product catalog, and begins taking advantage of the security features,
most of the details are handled by software.\\

By analogy, consider how you use an automated teller machine. A
complicated protocol has been defined by the banking industry to handle
transactions, but the mechanics of that protocol are hidden from you.
All you need to know is how to insert your card, type your personal
identification number, and request transactions. Ideally, a Secure
Mosaic tool will offer the same level of simplicity.\\

The Terisa Systems approach is by no means the only mechanism for
securing Mosaic or for conducting commerce on the Internet. A number of
firms have announced plans for Internet secure-transaction schemes.

\subsection{The clang of the cash register}
A New Hampshire startup, NetMarket, began offering online services in
the summer of 1994, selling products such as compact discs and flowers
over the Internet. It claims to have conducted the first secured Mosaic
transaction on August 10, 1994, using a secure version of X Mosaic that
allows users to place orders. NetMarket's secure Mosaic uses PGP as the
encryption code, which the company licensed from ViaCrypt.\\

Of course, relatively few end users are running X/Windows on Unix
workstations. ``PGP is the solution we're doing now,'' said Dan Kohn,
president of Net Market. ``It's technology we happen to like because
it's cheap for providers and free for consumers.'' Kohn is open to
adopting other standards more likely to be deployed for the masses as
they evolve. ``The best model would be for firms offering online
services such as ourselves to pay reasonable fees for server technology
and for secure clients to be free.''\\

Another company working on secure Web client technology is Netscape
Communications Corp., formerly known as Mosaic Communications Corp. The
company, whose principals include Jim Clark (the founder of Silicon
Graphics) and Marc Andreessen (the co-author of the original X Mosaic
from NCSA), is developing commercial Web client and server products.\\

Netscape's approach to securing Mosaic also relies on RSA Data
Security's public-key scheme, but its protocol enhancements differ from
those proposed by Terisa Systems. It proposes to implement a Secure
Socket Library to build a system it claims will be bulletproof.
Netscape has announced a partnership with First Data Card Services,
which is the world's largest credit-card authorization firm. Other
partners include three banks: First Interstate, Old Kent, and Norwest
Card Services.

\subsection{Money makes the world go round}
Even when new Web protocols are standardized and software is available,
there will be lag time before security-enhanced client software is
installed on millions of desktops. Not all companies are waiting for
these developments. Some are devising ways to provide a measure of
security today. One such firm is Open Market, which is building on
technology devised at the Massachusetts Institute of Technology.\\

Open Market provides a StoreBuilder toolkit for firms that want to
conduct business on the Internet. Its scheme works with existing Web
clients and uses a payment URL (Uniform Resource Locator) that encodes
price and other data in the information sent to the client program.\\

When a consumer selects an item, the URL is sent to a payment server,
which authenticates the user and processes the payment. The payment
server then redirects the client program to use an access URL. The
model supports ordering of information goods such as delivery of
magazine or database excerpts on demand as well as traditional hard
goods. Open Market's partners include the Lexis/Nexis service and
Digital Equipment Corp.\\

New security technologies and alliances are being announced at a
breathtaking pace. For example, Microsoft has announced an alliance
with Visa International. Microsoft recently announced its own online
information service (called Marvel) and purchased Intuit, the leading
maker of PC personal finance software.\\

Other new ventures such as First Virtual aim to provide transaction
services without a need for security-enhanced software (see Section
\ref{rush_commerce}). Purchases are authenticated off-line via a relatively
simple e-mail confirmation instead of more complex schemes.\\

During 1995, we probably will see many more alliances out to provide
secure Internet transaction technologies. A shakeout also is likely as
the market settles on technologies and players within and outside the
Internet context.\\

Advocates of commerce on the Internet maintain that the more
complicated schemes are essential for the Net to replace conventional
means of conducting business. To prevent a Tower of Babel of complex
new schemes, the newly formed W3 Organization is trying to foster
cooperation among technology vendors. Tim Berners-Lee, the chief
architect of the World-Wide Web, has convened meetings of major players
on behalf of the consortium.\\

Berners-Lee is optimistic that there will be some convergence. He also
points out that HTTP was designed under the assumption that multiple
schemes would coexist. According to Berners-Lee, the top layer of
Secure HTTP transactions will negotiate the level of security and the
exact scheme to be used. This will allow users and vendors to take
advantage of future encryption schemes that have not yet been invented.

